\chapter{Interpreting the Multiplier: Shadow Prices and the Envelope Theorem}

The multiplier $\lambda$ is more than algebraic clutter: at the optimum
it is the rate at which the optimised value $V$ responds to a relaxation of
the constraint. In consumer theory this is the marginal utility of income;
in producer theory the shadow price of capital or labour. The envelope
theorem makes this rigorous and supplies the tool for comparative statics.

\section{Core Rules}

\begin{rulebox}{Value function}
Consider the parametric problem
\[
V(b) \;=\; \max_{(x, y)} f(x, y) \quad \text{s.t.} \quad g(x, y) = b.
\]
$V$ is the \emph{value function} or indirect objective.
\end{rulebox}

\begin{rulebox}{Envelope theorem (one constraint)}
If $f, g$ are $C^1$ and the optimum $(x^*(b), y^*(b))$ varies smoothly in
$b$, then
\[
\dd{V}{b} \;=\; \lambda^*(b).
\]
Thus $\lambda^*$ is the \textbf{shadow price} of the constraint: it tells
you how much the optimised objective rises when the constraint is loosened
by one unit. In consumer problems $V(M)$ is the \emph{indirect utility} and
$\lambda^* = \partial V/\partial M$ is the marginal utility of income.
\end{rulebox}

\begin{rulebox}{Why ``envelope''?}
At the unconstrained level
$\mathcal{L}^*(b) = f(x^*, y^*) + \lambda^*(b - g(x^*, y^*)) = f(x^*, y^*)
= V(b)$ at the optimum. Differentiating $\mathcal{L}^*$ with respect to
$b$, the $x^*$ and $y^*$ contributions vanish thanks to the FOCs (chain
rule with $\mathcal{L}_x = \mathcal{L}_y = 0$), leaving only the explicit
$b$-dependence: $\partial \mathcal{L}/\partial b = \lambda^*$.
\end{rulebox}

\begin{rulebox}{General parameter}
For any parameter $\alpha$ entering $f$ or $g$, the envelope theorem gives
\[
\dd{V}{\alpha} \;=\; \pd{\mathcal{L}}{\alpha}\bigg|_{x^*, y^*, \lambda^*}.
\]
Only the \emph{explicit} dependence on $\alpha$ matters.
\end{rulebox}

\begin{rulebox}{Sign and units of $\lambda^*$}
For a maximisation problem with the convention
$\mathcal{L} = f + \lambda(b - g)$, $\lambda^* \geq 0$: relaxing a binding
constraint cannot hurt. Its units are (units of $f$) / (units of $g$),
``utils per euro'' in consumer theory, ``additional output per additional
input'' in production.
\end{rulebox}

\begin{rulebox}{Comparative statics (a brief preview)}
Given $V(\alpha)$, the dependence of $x^*(\alpha)$ on $\alpha$ can be
analysed via implicit differentiation of the FOC system. The envelope
theorem allows the derivative $\mathrm{d}V/\mathrm{d}\alpha$ to be computed
\emph{without} solving for $x^*(\alpha)$ explicitly.
\end{rulebox}

\begin{rulebox}{Duality: output maximisation and cost minimisation}
Consider the pair
\[
  \text{(P)}\ \max_{\mathbf{x}}F(\mathbf{x})\ \text{s.t.}\ \mathbf{w}\cdot\mathbf{x}=B,
  \qquad
  \text{(D)}\ \min_{\mathbf{x}}\mathbf{w}\cdot\mathbf{x}\ \text{s.t.}\ F(\mathbf{x})=\bar Q .
\]
Both have the same tangency conditions
$F_{x_i}/F_{x_j}=w_i/w_j$, hence the same \emph{expansion path}; they differ
only in which constraint fixes the point on it.
\begin{itemize}[leftmargin=1.6em,itemsep=0.3ex]
  \item If $\bar Q=Q^*$, the maximal output of (P), then the bundle
        $\mathbf{x}^*$ of (P) also solves (D), the minimal cost is exactly
        $B$, and the multipliers are reciprocal:
        $\mu_{\text{cost}}=1/\lambda_{\text{output}}$. ($\lambda$ is output
        per extra euro, $\mu$ is euros per extra unit of output: the
        marginal cost.)
  \item If $\bar Q\neq Q^*$, the answer moves along the expansion path. For
        $F$ homogeneous of degree $r$, scale the bundle by
        $t=(\bar Q/Q^*)^{1/r}$: inputs and cost are multiplied by $t$.
  \item The same duality holds in geometry: the shape with the largest
        volume for a given surface is the shape with the smallest surface
        for a given volume, and again the two multipliers are reciprocal.
\end{itemize}
In consumer theory the same pair is utility maximisation and expenditure
minimisation (Chapter~1).

\textbf{Monotone transformations.} Replacing the objective $f$ by $G\circ f$
with $G$ strictly increasing \emph{on the range of values $f$ takes on the
constraint set} leaves the optimiser unchanged (an ordinal property,
Chapter~3) and multiplies the multiplier by $G'(f^*)$. If $G$ is increasing
only on part of that range ($G(t)=t^{2}$ where $f$ can be negative), new
critical points and new optima may appear.
\end{rulebox}

\section{Worked Example}

\begin{worked}{$U=x^{0.4}y^{0.6}$, $p_{x}=2$, $p_{y}=3$, $M=100$: find the optimum, the multiplier, and verify $\lambda^{*}=\mathrm{d}V/\mathrm{d}M$}
\wstep{Step 1. Solve the problem.}
The Cobb--Douglas demands (Chapter~7) put a fixed expenditure share on each
good:
\[
  x^{*}=\frac{0.4M}{p_{x}}=\frac{40}{2}=20,
  \qquad
  y^{*}=\frac{0.6M}{p_{y}}=\frac{60}{3}=20 .
\]
Budget check: $2(20)+3(20)=40+60=100$. \checkmark

\wstep{Step 2. Evaluate the value function.}
\[
  V(M)=U(x^{*},y^{*})=20^{0.4}\cdot20^{0.6}=20^{1}=20 .
\]

\wstep{Step 3. Recover the multiplier from a first-order condition.}
$U_{x}=0.4x^{-0.6}y^{0.6}=0.4\bigl(y/x\bigr)^{0.6}$, which at
$x^{*}=y^{*}=20$ is $0.4$. From $U_{x}=\lambda p_{x}$,
\[
  \lambda^{*}=\frac{U_{x}}{p_{x}}=\frac{0.4}{2}=0.2 .
\]
Cross-check with the other condition: $U_{y}=0.6(x/y)^{0.4}=0.6$, and
$0.6/3=0.2$. \checkmark\ The two must agree --- that agreement \emph{is} the
tangency condition, and it is a free check on the whole solution.

\wstep{Step 4. Now verify the interpretation directly.}
For general $M$ the demands are $x^{*}=0.2M$ and $y^{*}=0.2M$, so
\[
  V(M)=(0.2M)^{0.4}(0.2M)^{0.6}=0.2M ,
  \qquad
  \dd{V}{M}=0.2=\lambda^{*} . \checkmark
\]
\[
  \boxed{\lambda^{*}=\pd{V}{M}:\ \text{the multiplier is the marginal value of the constraint.}}
\]

\wstep{Step 5. Use it as intended --- an approximation, without re-solving.}
If income rises to $M=105$, the predicted utility is
\[
  V(105)\approx V(100)+\lambda^{*}\cdot5=20+1=21 ,
\]
and here the prediction is exact ($0.2\times105=21$) because $V$ happens to
be linear in $M$. In general it is a first-order approximation, accurate for
small changes.
\wtrap{The approximation is local. For a large parameter change the
underlying solution may move to a different regime --- a constraint that was
slack becomes binding, or a variable hits zero --- and the multiplier
computed at the old optimum no longer applies. ``Small change'' is a real
qualification, not a politeness.}

\wstep{Step 6. Units, which are how you check any shadow price.}
$\lambda$ has the units of \emph{objective per unit of constraint}. Here:
utils per euro. In the cost-minimisation problem of Chapter~7 it was euros
per unit of output --- marginal cost. In a linear programme it is profit per
machine-hour. Getting the units right identifies what the number means, and
a $\lambda$ whose units make no sense signals a mis-specified constraint.
\wtrap{For a \emph{utility} problem, $\lambda$ has no absolute meaning. Utility
is ordinal: replacing $U$ by $U^{3}$ describes identical preferences,
leaves $x^{*},y^{*}$ unchanged, and multiplies $\lambda$ by
$3U^{2}=1200$. The multiplier is only meaningful when the objective is
measured in real units --- money, output, time --- which is why shadow
prices in cost, profit and programming problems are quoted and marginal
utilities of income are not.}

\wstep{Step 7. Where the ``envelope'' name comes from.}
$V(M)$ is the upper envelope of the family of curves $M\mapsto U(x,y)$
indexed by fixed bundles $(x,y)$: each fixed bundle's utility is a
horizontal line, affordable only once $M$ is large enough, and $V$ traces
the highest one available at each $M$. At every $M$ the envelope touches one
member of the family, and there the two have the same slope --- which is
exactly the theorem.

\wstep{Step 8. Sign, in general.}
Relaxing a binding constraint can never hurt, so for a maximisation problem
with the convention of Chapter~7, $\lambda^{*}\ge0$; and $\lambda^{*}=0$
precisely when the constraint is not binding. That last observation is the
whole content of complementary slackness in Chapter~10.
\end{worked}

\section{Practice Problems}

\begin{problems}
  \item In problem 1 of the previous chapter ($\max x y$ s.t.\
        $2 x + 3 y = 12$, with $\lambda^* = 1$, $V^* = 6$), verify the
        envelope theorem: change $12 \to 13$, recompute $V$, and check
        $\Delta V \approx \lambda^* \Delta b$.
  \item For $\max \ln x + \ln y$ s.t.\ $x + y = M$, derive $V(M)$ and
        verify $V'(M) = \lambda^*$.
  \item For Cobb--Douglas demand $x^* = \alpha M/p_x$,
        $y^* = (1 - \alpha) M/p_y$ from $U = x^\alpha y^{1 - \alpha}$,
        compute the indirect utility $V(p_x, p_y, M)$.
  \item For the indirect utility from problem 3, verify
        $\partial V/\partial M = \lambda^*$ (you may need to compute
        $\lambda^*$ from the FOC).
  \item Cost minimisation: $\min p_K K + p_L L$ s.t.\
        $K^{1/2} L^{1/2} = Q_0$. We found
        $C^*(Q_0) = 2 Q_0 \sqrt{p_K p_L}$. Compute
        $\mathrm{d}C^*/\mathrm{d}Q_0$ and identify the multiplier
        $\lambda^*$ (here the multiplier on the output constraint).
  \item In problem 9 of the previous chapter ($\min x^2 + y^2$ s.t.\
        $x + y = b$), find $V(b)$ and verify $V'(b) = \lambda^*(b)$.
  \item For utility $U = \alpha \ln x + \beta \ln y$ with budget
        $p_x x + p_y y = M$, derive the indirect utility $V$ and the
        marginal utility of income $\partial V/\partial M$.
  \item Maximise $U = x y$ on $p_x x + p_y y = M$ as a function of
        parameters: find $x^*, y^*, V(p_x, p_y, M)$. Then use the envelope
        theorem to compute $\partial V/\partial p_x$ without re-solving.
  \item A firm maximises $Q = K^{1/2} L^{1/2}$ subject to
        $r K + w L = TC$. We found $Q^* = TC/(2 \sqrt{r w})$. Use the
        envelope theorem to compute $\partial Q^*/\partial TC$ and
        $\partial Q^*/\partial r$ from $\mathcal{L}$ directly, and confirm
        with direct differentiation.
  \item Maximise $f(x, y) = 2 x + 3 y$ on $x^2 + y^2 = r^2$ for $r > 0$.
        Find $V(r) = f^*(r)$ and $\lambda^*(r)$. Verify $V'(r) = ?$:
        careful, the constraint is $g = r^2$, so the multiplier corresponds
        to changes in $r^2$, not $r$.
  \item In problem 1, suppose the budget changes from $12$ to
        $12 + \varepsilon$. Use $\lambda^*$ to estimate the new $V$ for
        $\varepsilon = 0.4$, then compute $V$ exactly and compare.
  \item A consumer has $U = x^{1/3} y^{2/3}$ and budget
        $p_x x + p_y y = M$ with $p_x = 1, p_y = 2, M = 30$. Find
        $x^*, y^*, \lambda^*$, and $V$. Then estimate how much $V$ changes
        if $M$ rises by $1$.
  \item For Cobb--Douglas with $\alpha = 2/5$, $p_x = 5, p_y = 4,
        M = 100$, find $x^*, y^*, \lambda^*$. By how much does $V$ change
        (approximately) when $p_x$ rises from $5$ to $5.1$? (Use
        envelope.)
  \item A firm maximises $Q = K^{1/3} L^{2/3}$ subject to
        $r K + w L = TC$. From problem 13 of the previous chapter,
        $K^* = TC/(3 r)$, $L^* = 2 TC/(3 w)$. Compute $Q^*(TC, r, w)$ and
        verify via the envelope theorem that
        $\partial Q^*/\partial w = -\lambda^* L^*$.
  \item In problem 17 of the previous chapter,
        $\min f = 3 x^2 + 2 y^2$ s.t.\ $x + y = b$ has
        $\lambda^*(b) = ?$. Find $V(b)$ explicitly and verify
        $V'(b) = \lambda^*(b)$.
\end{problems}

\section{Solutions}

\begin{solutions}
  \item Original: $\max x y$ on $2 x + 3 y = 12$: $V(12) = 6$,
        $\lambda^* = 1$. New: $2 x + 3 y = 13$. Resolve: $y = 2 \lambda$,
        $x = 3 \lambda$, constraint
        $6 \lambda + 6 \lambda = 12 \lambda = 13 \Rightarrow
        \lambda = 13/12$, $x = 13/4$, $y = 13/6$,
        $V(13) = 13/4 \cdot 13/6 = 169/24 \approx 7.042$. Then
        $\Delta V = 7.042 - 6 = 1.042$, while
        $\lambda^* \Delta b = 1 \cdot 1 = 1$. Close, as expected (the
        first-order approximation).

  \item FOCs: $1/x = \lambda$, $1/y = \lambda \Rightarrow x = y$.
        Constraint: $2 x = M \Rightarrow x^* = y^* = M/2$.
        $\lambda^* = 2/M$. $V(M) = 2 \ln(M/2) = 2 \ln M - 2 \ln 2$.
        $V'(M) = 2/M = \lambda^*$. \checkmark

  \item $V(p_x, p_y, M) = (x^*)^\alpha (y^*)^{1 - \alpha}
        = \left(\dfrac{\alpha M}{p_x}\right)^\alpha
          \left(\dfrac{(1 - \alpha) M}{p_y}\right)^{1 - \alpha}
        = \dfrac{\alpha^\alpha (1 - \alpha)^{1 - \alpha} M}{p_x^\alpha
          p_y^{1 - \alpha}}$. (Linear in $M$.)

  \item $\partial V/\partial M = \alpha^\alpha
        (1 - \alpha)^{1 - \alpha}/(p_x^\alpha p_y^{1 - \alpha})$. To check
        $\lambda^*$: from FOC,
        $\alpha x^{\alpha - 1} y^{1 - \alpha} = \lambda p_x \Rightarrow
        \lambda = \dfrac{\alpha U^*}{p_x x^*} = \dfrac{\alpha V}{\alpha M}
        = \dfrac{V}{M}$. Since $V$ is linear in $M$, $V/M = $ same
        constant $= \partial V/\partial M$. \checkmark

  \item $C^*(Q_0) = 2 Q_0 \sqrt{p_K p_L}$. Derivative
        $\mathrm{d}C^*/\mathrm{d}Q_0 = 2 \sqrt{p_K p_L}$. The constraint
        was $Q_0 - K^{1/2} L^{1/2} = 0$; in the minimisation
        $\mathcal{L} = C + \lambda(Q_0 - K^{1/2} L^{1/2})$, envelope gives
        $\partial C^*/\partial Q_0 = \lambda^*$. From FOC
        $p_K = \tfrac{\lambda}{2} K^{-1/2} L^{1/2}
        = \tfrac{\lambda}{2}(L^*/K^*)^{1/2}$. With $L^*/K^* = p_K/p_L$:
        $p_K = \tfrac{\lambda}{2} \sqrt{p_K/p_L} \Rightarrow
        \lambda = 2 p_K \sqrt{p_L/p_K} = 2 \sqrt{p_K p_L}$. Matches.

  \item FOCs: $2 x = \lambda$, $2 y = \lambda \Rightarrow x = y = b/2$.
        $V(b) = 2(b/2)^2 = b^2/2$. $\lambda^* = b$.
        $V'(b) = b = \lambda^*$. \checkmark

  \item Demands: $x^* = \dfrac{\alpha M}{(\alpha + \beta) p_x}$,
        $y^* = \dfrac{\beta M}{(\alpha + \beta) p_y}$.
        $V = \alpha \ln x^* + \beta \ln y^* = (\alpha + \beta) \ln M
        + \text{(terms in $\alpha, \beta, p_x, p_y$ not in $M$)}$. So
        $\partial V/\partial M = (\alpha + \beta)/M$. From FOC
        $\lambda^* = (\alpha + \beta)/M$ (problem 6 of the previous
        chapter). \checkmark

  \item FOCs: $y = \lambda p_x$, $x = \lambda p_y$. Divide:
        $y/x = p_x/p_y \Rightarrow p_x x = p_y y$. Budget:
        $2 p_x x = M \Rightarrow x^* = M/(2 p_x)$, $y^* = M/(2 p_y)$.
        $V = x^* y^* = M^2/(4 p_x p_y)$.
        $\lambda^* = y^*/p_x = M/(2 p_x p_y)$. Envelope:
        $\partial V/\partial p_x$ comes from
        $\partial \mathcal{L}/\partial p_x = -\lambda^* x^*$ (the only
        term in $\mathcal{L}$ depending explicitly on $p_x$ is
        $-\lambda p_x x$). Hence
        $\partial V/\partial p_x = -\lambda^* x^*
        = -\dfrac{M}{2 p_x p_y} \cdot \dfrac{M}{2 p_x}
        = -\dfrac{M^2}{4 p_x^2 p_y}$. Direct check:
        $\partial(M^2/(4 p_x p_y))/\partial p_x = -M^2/(4 p_x^2 p_y)$.
        \checkmark

  \item $\mathcal{L} = K^{1/2} L^{1/2} + \lambda(TC - r K - w L)$. By
        envelope, $\partial Q^*/\partial TC = \lambda^*$ and
        $\partial Q^*/\partial r = -\lambda^* K^*$. From problem 11 of the
        previous chapter: $K^* = TC/(2 r)$, $\lambda^* = $ ? From FOC
        $\tfrac{1}{2} K^{-1/2} L^{1/2} = \lambda r$, with $L^*/K^* = r/w$
        gives $\lambda^* = \tfrac{1}{2 r} \sqrt{r/w} = 1/(2 \sqrt{r w})$.
        Then $\partial Q^*/\partial TC = 1/(2 \sqrt{r w})$, matches
        direct differentiation of $Q^* = TC/(2 \sqrt{r w})$.
        $\partial Q^*/\partial r = -\dfrac{1}{2 \sqrt{r w}} \cdot
        \dfrac{TC}{2 r} = -\dfrac{TC}{4 r \sqrt{r w}}$. Direct:
        $Q^* = TC \cdot \tfrac{1}{2} r^{-1/2} w^{-1/2}$,
        $\partial/\partial r = TC \cdot \tfrac{1}{2} \cdot
        (-\tfrac{1}{2}) r^{-3/2} w^{-1/2}
        = -\dfrac{TC}{4 r^{3/2} \sqrt{w}}
        = -\dfrac{TC}{4 r \sqrt{r w}}$. \checkmark

  \item Constraint $g = x^2 + y^2 = r^2$,
        $\mathcal{L} = 2 x + 3 y + \lambda(r^2 - x^2 - y^2)$. From
        problem 7 of the previous chapter with $13 \to r^2$:
        $x = 2/(2 \lambda) = 1/\lambda$, $y = 3/(2 \lambda)$,
        $1/\lambda^2 + 9/(4 \lambda^2) = 13/(4 \lambda^2) = r^2
        \Rightarrow \lambda^* = \sqrt{13}/(2 r)$. Then
        $x^* = 2 r/\sqrt{13}$, $y^* = 3 r/\sqrt{13}$,
        $V(r) = 2(2 r/\sqrt{13}) + 3(3 r/\sqrt{13}) = 13 r/\sqrt{13}
        = r \sqrt{13}$. Envelope:
        $\partial V/\partial (r^2) = \lambda^*$. Indeed $V = r \sqrt{13}$
        so $\partial V/\partial r = \sqrt{13}$; and
        $\partial V/\partial(r^2) = \partial V/\partial r \cdot 1/(2 r)
        = \sqrt{13}/(2 r) = \lambda^*$. \checkmark

  \item $V(12.4) \approx V(12) + \lambda^* \cdot 0.4 = 6 + 0.4 = 6.4$.
        Exact: $V(b) = (b/(2 \cdot 2))(b/(3 \cdot 2)) \cdot ?$ Let's redo:
        FOCs $y = 2 \lambda$, $x = 3 \lambda$, $2 x + 3 y = b$,
        $6 \lambda + 6 \lambda = 12 \lambda = b$, $\lambda = b/12$,
        $x = b/4$, $y = b/6$, $V(b) = b^2/24$. So
        $V(12.4) = 153.76/24 \approx 6.407$. Approximation $6.4$ is
        accurate to two decimals.

  \item Cobb--Douglas with $\alpha = 1/3, \beta = 2/3$. Demands:
        $x^* = \alpha M/p_x = (1/3) \cdot 30 / 1 = 10$ and
        $y^* = \beta M / p_y = (2/3) \cdot 30 / 2 = 10$.
        $V = (x^*)^{1/3}(y^*)^{2/3} = 10^{1/3} \cdot 10^{2/3} = 10$. From
        the FOC $\alpha x^{\alpha - 1} y^{1 - \alpha} = \lambda p_x$,
        i.e.\ $(1/3) \cdot 10^{-2/3} \cdot 10^{2/3} = \lambda \cdot 1$, so
        $\lambda^* = 1/3$. Cross-check via direct differentiation:
        $V(M) = (\alpha M/p_x)^\alpha (\beta M/p_y)^\beta
        = M \cdot (1/3)^{1/3}(1/3)^{2/3} \cdot 1^{-1/3} 2^{-2/3}$, but
        more cleanly, plug in $p_x = 1, p_y = 2$:
        $V(M) = (M/3)^{1/3}(M/3)^{2/3} = M/3$, so
        $V'(M) = 1/3 = \lambda^*$. \checkmark A budget rise of $1$ raises
        $V$ by approximately $1/3$.

  \item Demands: $x^* = (2/5) \cdot 100/5 = 8$,
        $y^* = (3/5) \cdot 100/4 = 15$. From FOC
        $\alpha/x = \lambda p_x$,
        $\lambda^* = (2/5)/(5 \cdot 8) = (2/5)/40 = 1/100$. By the
        envelope theorem,
        $\partial V/\partial p_x = -\lambda^* x^* = -(1/100) \cdot 8
        = -0.08$. So $\Delta V \approx -0.08 \cdot 0.1 = -0.008$: a small
        loss in utility as $p_x$ rises.

  \item $Q^* = (K^*)^{1/3}(L^*)^{2/3}
        = (TC/(3 r))^{1/3}(2 TC/(3 w))^{2/3}$. Factoring out $TC$,
        \[
          Q^* = TC \cdot (1/(3 r))^{1/3}(2/(3 w))^{2/3}
              = TC \cdot 2^{2/3}/(3 r^{1/3} w^{2/3}).
        \]
        Envelope theorem says
        $\partial Q^*/\partial w = -\lambda^* L^*$. Direct:
        $\partial Q^*/\partial w = TC \cdot 2^{2/3}/3 \cdot r^{-1/3}
        \cdot (-2/3) w^{-5/3}$. From FOC
        $\tfrac{1}{3} K^{-2/3} L^{2/3} = \lambda r$, with $L/K = 2 r/w$,
        we get $\lambda = \tfrac{1}{3 r}(L/K)^{2/3}
        = \tfrac{1}{3 r}(2 r/w)^{2/3}$, so
        $\lambda^* L^* = \tfrac{1}{3 r}(2 r/w)^{2/3} \cdot 2 TC/(3 w)
        = \tfrac{2 TC}{9 r w}(2 r/w)^{2/3}
        = \tfrac{2 TC \cdot 2^{2/3} r^{2/3}}{9 r w \cdot w^{2/3}}
        = \tfrac{2 \cdot 2^{2/3} TC}{9 r^{1/3} w^{5/3}}$. Direct:
        $\tfrac{2^{2/3} TC \cdot (-2)}{9 r^{1/3} w^{5/3}}
        = -\tfrac{2 \cdot 2^{2/3} TC}{9 r^{1/3} w^{5/3}}$. Matches
        $-\lambda^* L^*$. \checkmark

  \item FOCs: $6 x = \lambda, 4 y = \lambda \Rightarrow x = \lambda/6,
        y = \lambda/4$. Constraint:
        $\lambda/6 + \lambda/4 = (2 \lambda + 3 \lambda)/12
        = 5 \lambda/12 = b \Rightarrow \lambda^*(b) = 12 b/5$. Then
        $x^* = 2 b/5, y^* = 3 b/5$.
        $V(b) = 3(2 b/5)^2 + 2(3 b/5)^2 = 12 b^2/25 + 18 b^2/25
        = 30 b^2/25 = 6 b^2/5$. $V'(b) = 12 b/5 = \lambda^*(b)$.
        \checkmark At $b = 5$: $\lambda^* = 12$, $V = 30$.
\end{solutions}

\section{Additional Problems}

\begin{problems}
  \item For $\max U = 3 x y$ s.t.\ $4 x + 5 y = b$, derive $V(b)$ and
        $\lambda^*(b)$. Verify $V'(b) = \lambda^*$.
  \item For $\max U = \ln x + 3 \ln y$ s.t.\ $2 x + y = b$, derive
        $V(b)$. Use $b = 20$ to estimate the effect of a one-unit budget
        increase.
  \item Cobb--Douglas with $\alpha = 1/4$, $p_x = 2, p_y = 3, M = 60$:
        find $x^*, y^*, \lambda^*$. By how much does $V$ change
        approximately if $p_y$ rises from $3$ to $3.1$?
  \item Cost minimisation $\min p_K K + p_L L$ s.t.\
        $K^{1/3} L^{2/3} = Q_0$. Find $C^*(Q_0, p_K, p_L)$ and verify
        $\partial C^*/\partial Q_0 = \lambda^*$.
  \item $\min f = x^2 + 2 y^2$ s.t.\ $2 x + y = b$: find $V(b)$
        explicitly and verify $V'(b) = \lambda^*(b)$.
  \item A firm maximises $Q = K^{2/3} L^{1/3}$ subject to
        $r K + w L = TC$. From the previous chapter, $K^* = 2 TC/(3 r)$,
        $L^* = TC/(3 w)$. Use the envelope theorem to compute
        $\partial Q^*/\partial TC$.
  \item For $\max U = x^{1/2} y^{1/2}$ on $p_x x + p_y y = M$ at
        $p_x = 4, p_y = 1, M = 40$, find $x^*, y^*, V, \lambda^*$. Use
        envelope to estimate $\Delta V$ when $p_x$ drops from $4$ to
        $3.9$.
  \item For $\max U = \ln(x + 1) + \ln(y + 1)$ s.t.\ $x + y = b$, find
        $V(b)$ and $\lambda^*(b)$ for general $b$. Check
        $V'(b) = \lambda^*$.
  \item For $\min f = (x - 2)^2 + (y - 3)^2$ s.t.\ $x + 2 y = b$, find
        $V(b)$ and $\lambda^*(b)$ for $b$ near $8$ (note: at $b = 8$ the
        constraint just touches the unconstrained min).
  \item For Cobb--Douglas $U = x^{1/2} y^{1/2}$ with
        $p_x = 1, p_y = 1, M = 10$, suppose a per-unit tax raises $p_x$
        from $1$ to $1.5$. Use envelope for a first-order estimate of
        $\Delta V$, then compute exact $\Delta V$ and compare.

  \item (Duality, constant returns.) A firm with $Q=K^{1/2}L^{1/2}$,
        $r=4$, $w=1$ maximises output with the budget $4K+L=80$.
        (a) Solve and find $\lambda^*$. (b) \emph{Without solving}, state
        the solution of $\min 4K+L$ s.t.\ $Q=20$, including its multiplier.
        (c) Still without solving, give the cost-minimising bundle for
        $Q=30$. (d) Confirm (b) and (c) by Lagrange.
  \item (Duality, decreasing returns.) $Q=K^{1/3}L^{1/3}$, $r=2$, $w=1$,
        budget $B=60$. (a) Find $\mathbf{x}^*$, $Q^*$ and $\lambda^*$.
        (b) Answer without solving: the minimal cost of producing $Q^*$, and
        its multiplier. (c) How do inputs and cost change if the target is
        $2Q^*$? Confirm with the cost function.
  \item (Geometric dual pair.) A closed box has a square base of side $a$
        and height $h$. (a) Maximise its volume for a surface area of $24$.
        (b) Minimise its surface for a volume of $8$. (c) Show that both
        give the same shape and that the two multipliers are reciprocal.
  \item (Cylinder.) A closed cylinder has radius $r$ and height $h$.
        (a) Maximise the volume for a surface of $54\pi$. (b) State, without
        solving, the solution of minimising the surface for the volume found
        in (a), then confirm by Lagrange and check $\lambda\mu=1$.
  \item (Monotone transformations.) On the constraint $x+y=10$ with
        $x,y>0$, compare the maximisers of $f=xy$, $\sqrt{f}$, $f^{2}$,
        $\mathrm{e}^{f}$ and $\ln f$, and the multipliers. Which of the five
        objectives are concave on $\R^{2}_{++}$? What changes for $f^{2}$
        if $x$ and $y$ may be negative?
  \item (Three inputs.) Maximise $Q=x_1^{0.2}x_2^{0.3}x_3^{0.5}$ subject
        to $2x_1+3x_2+5x_3=100$. Use the expenditure-share rule, then check
        the FOCs and interpret $\lambda^*$.
  \item Maximise $Q=(x_1x_2x_3)^{1/4}$ subject to $x_1+2x_2+4x_3=48$.
        Find $\mathbf{x}^*$, $Q^*$ and $\lambda^*$, and verify
        $\lambda^*=\tfrac34Q^*/B$.
\end{problems}

\section{Additional Solutions}

\begin{solutions}
  \item FOCs: $3 y = 4 \lambda$, $3 x = 5 \lambda$, so
        $y = 4 \lambda/3$, $x = 5 \lambda/3$. Constraint:
        $4 \cdot 5 \lambda/3 + 5 \cdot 4 \lambda/3 = 40 \lambda/3 = b
        \Rightarrow \lambda^*(b) = 3 b/40$. Then $x^* = b/8$,
        $y^* = b/10$. $V(b) = 3 \cdot (b/8)(b/10) = 3 b^2/80$.
        $V'(b) = 6 b/80 = 3 b/40 = \lambda^*(b)$. \checkmark At
        $b = 40$: $V = 60$, $\lambda^* = 3$.

  \item FOCs: $1/x = 2 \lambda$,
        $3/y = \lambda \Rightarrow y = 3/\lambda = 6 x$. Constraint:
        $2 x + 6 x = 8 x = b \Rightarrow x^* = b/8$, $y^* = 3 b/4$.
        $\lambda^*(b) = 1/(2 x^*) = 4/b$.
        $V(b) = \ln(b/8) + 3 \ln(3 b/4) = 4 \ln b + \text{const}$.
        $V'(b) = 4/b = \lambda^*$. \checkmark At $b = 20$,
        $\lambda^* = 1/5$, so a one-unit budget increase raises $V$ by
        approximately $0.2$.

  \item Cobb--Douglas demands: $x^* = \alpha M/p_x = (1/4)(60)/2
        = 15/2$; $y^* = (3/4)(60)/3 = 15$. From FOC
        $\alpha/x = \lambda p_x$ for the form
        $U = \alpha \ln x + (1 - \alpha) \ln y$, but here
        $U = x^\alpha y^{1 - \alpha}$. Use
        $\lambda^* = V/M \cdot (\alpha + (1-\alpha))/1$. Actually, for
        Cobb--Douglas $V$ is linear in $M$:
        $V(M) = \alpha^\alpha (1 - \alpha)^{1 - \alpha} M /
        (p_x^\alpha p_y^{1 - \alpha})$. At our parameters
        $V = (1/4)^{1/4}(3/4)^{3/4} \cdot 60 / (2^{1/4} \cdot 3^{3/4})$.
        Then $\lambda^* = V/M$. Envelope:
        $\partial V/\partial p_y = -\lambda^* y^* = -(V/60) \cdot 15
        = -V/4$. So $\Delta V \approx -(V/4) \cdot 0.1 = -V/40$.
        Numerically: $(1/4)^{1/4} \approx 0.707$,
        $(3/4)^{3/4} \approx 0.806$, $2^{1/4} \approx 1.189$,
        $3^{3/4} \approx 2.280$.
        $V \approx 0.707 \cdot 0.806 \cdot 60/(1.189 \cdot 2.280)
        \approx 34.2/2.71 \approx 12.6$. Then
        $\Delta V \approx -12.6/40 \approx -0.315$.

  \item Cost min: from FOC
        $w_1 - (\lambda/3) x_1^{-2/3} x_2^{2/3} = 0$ and
        $w_2 - (2\lambda/3) x_1^{1/3} x_2^{-1/3} = 0$ (with
        $w_1 = p_K, w_2 = p_L$, $x_1 = K, x_2 = L$). Divide:
        $p_K/p_L = (1/2)(x_2/x_1) \Rightarrow x_2 = 2 (p_K/p_L) x_1$,
        i.e.\ $L = 2 (p_K/p_L) K$. Constraint $K^{1/3} L^{2/3} = Q_0$:
        $K^{1/3}(2 p_K/p_L)^{2/3} K^{2/3} = Q_0 \Rightarrow
        K(2 p_K/p_L)^{2/3} = Q_0 \Rightarrow
        K^* = Q_0 (p_L/(2 p_K))^{2/3}$ and
        $L^* = 2 (p_K/p_L) K^* = Q_0 \cdot 2^{1/3} (p_K/p_L)^{1/3}$.
        $C^* = p_K K^* + p_L L^*
        = Q_0 [p_K (p_L/(2 p_K))^{2/3}
        + p_L \cdot 2^{1/3}(p_K/p_L)^{1/3}]
        = Q_0 \cdot p_K^{1/3} p_L^{2/3}[2^{-2/3} + 2^{1/3}]
        = Q_0 \cdot p_K^{1/3} p_L^{2/3} \cdot 3 \cdot 2^{-2/3}$ (since
        $2^{-2/3} + 2^{1/3} = 2^{-2/3}(1 + 2) = 3 \cdot 2^{-2/3}$). So
        $C^*(Q_0) = 3 \cdot 2^{-2/3} Q_0 \, p_K^{1/3} p_L^{2/3}$.
        $\mathrm{d}C^*/\mathrm{d}Q_0 = 3 \cdot 2^{-2/3} p_K^{1/3}
        p_L^{2/3} = \lambda^*$. \checkmark

  \item FOCs: $2 x = 2 \lambda \Rightarrow x = \lambda$;
        $4 y = \lambda \Rightarrow y = \lambda/4$. Constraint:
        $2 \lambda + \lambda/4 = 9 \lambda/4 = b \Rightarrow
        \lambda^*(b) = 4 b/9$. Then $x^* = 4 b/9$, $y^* = b/9$.
        $V(b) = (4 b/9)^2 + 2(b/9)^2 = 16 b^2/81 + 2 b^2/81
        = 18 b^2/81 = 2 b^2/9$. $V'(b) = 4 b/9 = \lambda^*(b)$.
        \checkmark

  \item $\mathcal{L} = K^{2/3} L^{1/3} + \lambda(TC - r K - w L)$. From
        FOC $(2/3) K^{-1/3} L^{1/3} = \lambda r$. At
        $K^* = 2 TC/(3 r)$, $L^* = TC/(3 w)$ we get $L^*/K^* = r/(2 w)$, so
        $\lambda^* = (2/(3 r))(L^*/K^*)^{1/3} = (2/(3 r))(r/(2 w))^{1/3}$.
        Envelope:
        \[
          \partial Q^*/\partial TC = \lambda^*
            = (2/(3 r))(r/(2 w))^{1/3} = 2^{2/3}/(3 r^{2/3} w^{1/3}).
        \]
        Direct: $Q^* = (K^*)^{2/3}(L^*)^{1/3}
        = (2 TC/(3 r))^{2/3}(TC/(3 w))^{1/3}
        = TC \cdot 2^{2/3}/(3 r^{2/3} w^{1/3})$, so
        $\partial Q^*/\partial TC = 2^{2/3}/(3 r^{2/3} w^{1/3})$.
        Matches. \checkmark

  \item Cobb--Douglas $\alpha = 1/2$: $x^* = M/(2 p_x) = 40/8 = 5$;
        $y^* = M/(2 p_y) = 20$. $V = \sqrt{x^* y^*} = \sqrt{100} = 10$.
        From additional problem 2 of the previous chapter,
        $\lambda^* = 1/(2 \sqrt{p_x p_y}) = 1/(2 \cdot 2) = 1/4$.
        Envelope: $\partial V/\partial p_x = -\lambda^* x^*
        = -(1/4)(5) = -5/4$. $\Delta V \approx (-5/4)(-0.1) = 0.125$.
        (When $p_x$ falls, utility rises.)

  \item FOCs: $1/(x + 1) = \lambda$,
        $1/(y + 1) = \lambda \Rightarrow x = y$. Constraint:
        $2 x = b \Rightarrow x^* = y^* = b/2$.
        $\lambda^*(b) = 1/(b/2 + 1) = 2/(b + 2)$.
        $V(b) = 2 \ln(b/2 + 1) = 2 \ln((b + 2)/2)$.
        $V'(b) = 2 \cdot (1/(b + 2)) = 2/(b + 2) = \lambda^*$.
        \checkmark

  \item FOCs: $2(x - 2) = \lambda$,
        $2(y - 3) = 2 \lambda \Rightarrow y - 3 = \lambda$. So
        $y - 3 = 2(x - 2) \Rightarrow y = 2 x - 1$. Constraint:
        $x + 2(2 x - 1) = 5 x - 2 = b \Rightarrow x^* = (b + 2)/5$,
        $y^* = 2(b + 2)/5 - 1 = (2 b - 1)/5$.
        $\lambda^*(b) = 2(x^* - 2) = 2((b + 2)/5 - 2) = 2(b - 8)/5$. At
        $b = 8$: $\lambda^* = 0$, $(x^*, y^*) = (2, 3)$, which is the
        unconstrained minimum (constraint touches it).
        $V(b) = (x^* - 2)^2 + (y^* - 3)^2
        = ((b - 8)/5)^2 + (2(b - 8)/5)^2
        = (b - 8)^2/25 + 4(b - 8)^2/25 = (b - 8)^2/5$.
        $V'(b) = 2(b - 8)/5 = \lambda^*(b)$. \checkmark

  \item Cobb--Douglas $\alpha = 1/2$, $p_x = 1, p_y = 1, M = 10$:
        $x^* = 5, y^* = 5$, $V = 5$,
        $\lambda^* = 1/(2 \sqrt{1 \cdot 1}) = 1/2$. Envelope:
        $\partial V/\partial p_x = -\lambda^* x^* = -5/2$. Approx
        $\Delta V \approx -2.5 \cdot 0.5 = -1.25$. Exact: at new prices
        $V_{\text{new}} = M/(2 \sqrt{p_x p_y}) = 10/(2 \sqrt{1.5})
        = 5/\sqrt{1.5} \approx 4.082$. Exact
        $\Delta V = 4.082 - 5 = -0.918$. The approximation overshoots in
        magnitude because $p_x$ changes by a large $50\%$; first-order
        accurate only for small changes.

  \item (a) Exponents $\tfrac12,\tfrac12$: equal cost shares, $4K=40$,
        $L=40$, so $K^*=10$, $L^*=40$, $Q^*=\sqrt{400}=20$.
        $\lambda^*=\dfrac{Q_K}{r}=\dfrac{\tfrac12\sqrt{L/K}}{4}=\dfrac{1}{4}$
        (also $Q^*/B=\tfrac{20}{80}$, constant returns).
        (b) The target equals $Q^*$, so the same bundle $(10,40)$ is cost
        minimising, the minimal cost is the budget $80$, and
        $\mu^*=1/\lambda^*=4$, the marginal cost of output.
        (c) Degree of homogeneity $1$: scale by $t=\tfrac{30}{20}=1.5$, giving
        $(15,60)$ at cost $120$.
        (d) $\mathcal{L}=4K+L+\mu(\bar Q-\sqrt{KL})$:
        $4=\tfrac{\mu}{2}\sqrt{L/K}$, $1=\tfrac{\mu}{2}\sqrt{K/L}$; dividing,
        $L=4K$, and $\sqrt{4K^{2}}=2K=\bar Q$. For $\bar Q=20$: $(10,40)$,
        and $1=\tfrac{\mu}{2}\cdot\tfrac12$ gives $\mu=4$. For $\bar Q=30$:
        $(15,60)$, cost $120$, $\mu=4$ again (constant marginal cost).
  \item (a) Equal exponents, equal spending: $2K=30$, $L=30$, so
        $K^*=15$, $L^*=30$, $Q^*=450^{1/3}\approx7.663$, and
        $\lambda^*=\dfrac{Q_L}{w}=\dfrac{Q^*}{3L^*}\approx0.0851$.
        (b) Minimal cost of $Q^*$: $60$, with $\mu^*=1/\lambda^*\approx11.74$.
        (c) $F$ is homogeneous of degree $\tfrac23$ (not to be confused with
        the input price $r=2$), so doubling output
        needs $t=2^{3/2}\approx2.828$: $K\approx42.43$, $L\approx84.85$, cost
        $60\cdot2^{3/2}\approx169.7$. Check: in general $K=B/4$, $L=B/2$, so
        $Q^*(B)=(B^{2}/8)^{1/3}=\tfrac12B^{2/3}$; inverting,
        $C(Q)=(2Q)^{3/2}$, and $C(2Q^*)=(4Q^*)^{3/2}\approx169.7$. Also
        $C'(Q^*)=3\sqrt{2Q^*}\approx11.74=\mu^*$. \checkmark
  \item (a) $\mathcal{L}=a^{2}h+\lambda(24-2a^{2}-4ah)$:
        $2ah=\lambda(4a+4h)$, $a^{2}=4\lambda a$, so $\lambda=\tfrac a4$ and
        $2ah=a^{2}+ah$, i.e.\ $h=a$. Then $6a^{2}=24$: $a=h=2$, volume $8$,
        $\lambda^*=\tfrac12$.
        (b) $\mathcal{L}=2a^{2}+4ah+\mu(8-a^{2}h)$: $4a+4h=2\mu ah$,
        $4a=\mu a^{2}$, so $\mu=\tfrac4a$ and $4a+4h=8h$, $h=a$. Then
        $a^{3}=8$: $a=h=2$, surface $24$, $\mu^*=2$.
        (c) Both give the cube of side $2$, and $\lambda^*\mu^*=1$: $\lambda$
        is extra volume per unit of surface, $\mu$ extra surface per unit of
        volume.
  \item (a) $\mathcal{L}=\pi r^{2}h+\lambda(54\pi-2\pi r^{2}-2\pi rh)$:
        $2\pi rh=\lambda(4\pi r+2\pi h)$ and $\pi r^{2}=2\pi\lambda r$, so
        $\lambda=\tfrac r2$ and $2\pi rh=2\pi r^{2}+\pi rh$, i.e.\ $h=2r$
        (height equals diameter). Then $6\pi r^{2}=54\pi$: $r=3$, $h=6$,
        $V=54\pi$, $\lambda^*=\tfrac32$.
        (b) By duality the same can, $r=3$, $h=6$, surface $54\pi$, with
        $\mu^*=\tfrac23$. Lagrange:
        $\mathcal{L}=2\pi r^{2}+2\pi rh+\mu(54\pi-\pi r^{2}h)$:
        $\mathcal{L}_h=2\pi r-\mu\pi r^{2}=0$ gives $\mu=\tfrac2r$;
        $\mathcal{L}_r=4\pi r+2\pi h-2\mu\pi rh=0$ gives $4r+2h=4h$, $h=2r$;
        then $2\pi r^{3}=54\pi$, $r=3$, $\mu=\tfrac23$, and
        $\lambda\mu=\tfrac32\cdot\tfrac23=1$.
  \item On the constraint with $x,y>0$, $f=xy$ takes values in $(0,25]$,
        and all four transforms $\sqrt t$, $t^{2}$, $\mathrm{e}^{t}$,
        $\ln t$ are strictly increasing on $(0,\infty)$: every objective has
        the same maximiser $(5,5)$. For $G\circ f$ the FOC reads
        $G'(xy)\,y=\lambda=G'(xy)\,x$, so $x=y=5$ and
        $\lambda=5G'(25)$: $\lambda=5$ for $f$, $\tfrac12$ for $\sqrt f$,
        $250$ for $f^{2}$, $5\mathrm{e}^{25}$ for $\mathrm{e}^{f}$, $\tfrac15$
        for $\ln f$. Concave on $\R^{2}_{++}$: $\sqrt{xy}$ (Cobb--Douglas of
        degree $1$) and $\ln f=\ln x+\ln y$; not $xy$, $(xy)^{2}$ or
        $\mathrm{e}^{xy}$ (their Hessians are indefinite, e.g.\ along $y=x$
        they grow like $t^{2}$, $t^{4}$, $\mathrm{e}^{t^{2}}$). If $x,y$ may
        be negative, $xy=x(10-x)$ takes every value $\le25$, and $t^{2}$ is
        decreasing for $t<0$. The FOC for $(xy)^{2}$, $2xy\cdot y=\lambda
        =2xy\cdot x$, now also holds at $(0,10)$ and $(10,0)$, where
        $(xy)^{2}=0$ is a minimum, and $(xy)^{2}=x^{2}(10-x)^{2}\to\infty$ as
        $x\to\pm\infty$: no maximum exists at all.
  \item Exponents sum to $1$; the expenditure on $x_i$ is its exponent
        times $B$: $2x_1=20$, $3x_2=30$, $5x_3=50$, so
        $\mathbf{x}^*=(10,10,10)$ and $Q^*=10^{0.2+0.3+0.5}=10$. FOCs
        $Q_{x_i}=\lambda w_i$: $0.2\cdot\tfrac{Q}{x_1}=0.2=2\lambda$,
        $0.3=3\lambda$, $0.5=5\lambda$, all with $\lambda^*=0.1$. One more
        euro of budget buys $0.1$ units of output; with constant returns
        $\lambda^*=Q^*/B$.
  \item Equal exponents $\tfrac14$: equal spending $16$ on each input,
        $\mathbf{x}^*=(16,8,4)$, $Q^*=(512)^{1/4}\approx4.757$. FOC:
        $\tfrac14\tfrac{Q}{x_1}=\lambda\cdot1$ gives
        $\lambda^*=\tfrac{Q^*}{64}\approx0.0743$ (the same from $x_2$, $x_3$:
        $\tfrac{Q}{4\cdot8\cdot2}$, $\tfrac{Q}{4\cdot4\cdot4}$). With
        $B=48$: $\tfrac34\cdot\tfrac{Q^*}{48}=\tfrac{Q^*}{64}$. \checkmark\
        (Homogeneity of degree $\tfrac34$ makes $Q^*(B)\propto B^{3/4}$, so
        $\mathrm{d}Q^*/\mathrm{d}B=\tfrac34Q^*/B$.)
\end{solutions}
